\documentclass[11pt,letterpaper]{article}
\usepackage[T1]{fontenc}
\usepackage[utf8]{inputenc}
\usepackage{lmodern,microtype}
\usepackage[margin=0.94in]{geometry}
\usepackage{amsmath,amssymb,amsthm,mathtools}
\usepackage{booktabs,longtable,array,tabularx}
\usepackage{xcolor,xurl,hyperref,enumitem,fancyhdr,needspace}
\definecolor{ink}{HTML}{18374D}
\hypersetup{colorlinks=true,linkcolor=ink,urlcolor=ink,citecolor=ink,
 pdftitle={PA194N: a stack-preserving modus-ponens machine for Con(PA)},
 pdfauthor={ChatGPT}}
\urlstyle{tt}
\setlength{\emergencystretch}{3em}
\setlength{\parindent}{0pt}
\setlength{\parskip}{4pt}
\setlist{itemsep=3pt,topsep=5pt}
\pagestyle{fancy}\fancyhf{}
\fancyhead[L]{\small PA194N: a PA consistency machine}
\fancyhead[R]{\small Construction, comparison and proof}
\fancyfoot[C]{\thepage}
\setlength{\headheight}{14pt}
\newcommand{\PA}{\mathrm{PA}}
\newcommand{\Con}{\operatorname{Con}}
\newcommand{\Prf}{\operatorname{Prf}}
\newcommand{\Halt}{\operatorname{Halt}}
\newcommand{\file}[1]{\texttt{\nolinkurl{#1}}}
\newcommand{\MN}{M_{194N}}
\newcommand{\AX}{\mathsf{AX}}
\newcommand{\MP}{\mathsf{MP}}
\newcommand{\SHIFT}{\mathsf S}
\newcommand{\POP}{\mathsf{Pop}}
\newcommand{\PUSH}{\mathsf{Push}}
\newcommand{\DISCARD}{\mathsf{Discard}}
\newtheorem{theorem}{Theorem}[section]
\newtheorem{lemma}[theorem]{Lemma}
\newtheorem{proposition}[theorem]{Proposition}
\newtheorem{corollary}[theorem]{Corollary}
\title{\textcolor{ink}{PA194N}\\A stack-preserving modus-ponens machine\\for the consistency of Peano arithmetic}
\author{Comparative research and executable verification}
\date{13 November 2026, Pacific time}
\begin{document}
\maketitle
\begin{abstract}
We compare the 41 research packets represented in the requested repository
snapshot and construct a new deterministic, binary, one-tape, blank-start
machine with \textbf{194 working states}. We do not obtain a verified
machine with fewer than 194 states. The accepted result is the requested
alternative: PA194N changes the successful modus-ponens command so that
it preserves the proof stack, while retaining the original Shift test
$t+h=1$. Its exact transition graph is absent from all eight previously
represented 194-state classes, including supplementary and rejected tables.

The construction retains the primary 67-reader, seven-fold axiom decoder.
One native transition bypasses a 19-instruction dummy-zero push. The resulting
source has 3,654 primitive sites and 4,755 symbolic outcomes. Both physical
checkers pass every outcome for arbitrary natural register values and an
arbitrary nonempty spare list; an independent checker rederives the generic
register-loop lemmas from the actual rows. Exact blank initialization takes
58,154 transitions and is checked against the complete tape.

A ten-entry command compiler proves completeness of the changed interface.
In particular, when the offset is zero, $\MP(0)$ alone resets an arbitrary current formula without
changing the stack. Together with syntactic proof annotation for soundness,
positive physical simulation, and the explicitly inherited PA calibration
of the separate-diagonal-head calculus, the ordinary mathematical argument gives
\[
 \PA\vdash\bigl(\neg\Halt(\MN,\mathrm{blank})
                  \leftrightarrow\Con(\PA)\bigr).
\]
This is not a newly kernel-checked Lean or Rocq endpoint. The archive supplies
the literal table, deterministic construction, independent replay, mathematical
proofs, complete corpus census, search records, and all LaTeX sources.
\end{abstract}
\clearpage
{\small\setlength{\parskip}{0pt}\tableofcontents}
\newpage

\section{Outcome, conventions, and exact identity}
\label{sec:result}

The primary deliverable is \file{construction/machine/pa194N.tm}.
The machine has working states $0,\ldots,193$, initial state $0$, and a
separate terminal label $194$. The head starts at coordinate zero of a
two-way infinite all-zero tape. The alphabet is $\{0,1\}$; every transition
writes one bit and moves one cell left or right. There are 388 rows and no
row for the halt label. No input, external initialization, stationary move,
or change of state-counting convention is used.

The exact dense table SHA-256 is
\begin{quote}\small
\nolinkurl{37f7e909724cab7fd667db28f89b9c81653d2576c033750e783b40b3289c04f9}.
\end{quote}
The native presentation starts in 74 and halts in 254. Its SHA-256 is
\begin{quote}\small
\nolinkurl{b38dfe8cbd6e84aadbea5912e20f25b0bbda8f3fad4fbec0c779a77c857a8dc8}.
\end{quote}
A supplied bijection checks every native/dense transition. All 194 working
states are reachable in the designated finite transition graph. This does
not assert that every state occurs on the actual blank run.

\begin{center}
\begin{tabular}{@{}lrrr@{}}
\toprule
Machine & Working states & Primitive source sites & Primitive outcomes\\
\midrule
Primary PA194 & 194 & 3,673 & 4,780\\
PA194C & 194 & 3,649 & 4,749\\
PA194CS & 194 & 3,646 & 4,745\\
\textbf{PA194N} & \textbf{194} & \textbf{3,654} & \textbf{4,755}\\
\bottomrule
\end{tabular}
\end{center}
PA194N offers a different command protocol; it is not a claim to have the
shortest source among the 194-state machines. The search did not establish
193 states or a lower bound excluding them.

\begin{theorem}[Ordinary PA calibration of the delivered table]
Let $H_N$ be the arithmetical assertion of a finite halting computation of
the specified literal machine from its blank start. Using the published
syntactic PA calibration of the retained numerical calculus,
\[
 \PA\vdash H_N\leftrightarrow
   \exists p\,\Prf_{\PA}(p,\ulcorner0=1\urcorner),
 \qquad
 \PA\vdash\neg H_N\leftrightarrow\Con(\PA).
\]
\end{theorem}
The proof has three distinct interfaces: the numerical-calculus/PA proof
translations, the new proof-to-command compiler and annotated command
execution, and the literal source/tape simulation. The latter two are
specified below for PA194N. The inherited logical theorem is described
explicitly in Section~\ref{sec:logic}; a separate conditional iPRA assembly
in the repository is not used as an unconditional proof of this theorem.
Actual nonhalting follows under the assumption that PA is consistent.

\section{The substantive change: do not push a dummy zero}
\label{sec:change}

Write a command store as $(t,w,d,h)$: the current formula, a numerical
Cantor stack, an offset, and the cached formula. Put
\[
 P(a,b)=\frac{(a+b)(a+b+1)}2+b,
 \quad [a;w]=P(a,w),\quad I(a,b)=4P(a,b).
\]
The ordinary zero-stack convention is literal: $[0;0]=0$. Formula code
zero denotes truth and code one denotes the designated contradiction.
For a modus-ponens parameter $x$, set $b=d+x$. The successful comparison is
$I(t,b)=h$. The three relevant interfaces are
\begin{center}
\begin{tabular}{@{}lll@{}}
\toprule
Controller & Successful $\MP(x)$ result & Shift halt test\\
\midrule
Original PA194 & $(b,[0;w],b,0)$ & $t+h=1$\\
Cached PA194C & $(0,w,0,b)$ & $h=1$\\
\textbf{PA194N} & $\boldsymbol{(b,w,b,0)}$ & $\boldsymbol{t+h=1}$\\
\bottomrule
\end{tabular}
\end{center}
The failed comparison in PA194N returns $(0,w,b,0)$. Its nonhalting Shift
is the original
\[
 \SHIFT(t,[a;v],d,h)=(h,v,d,a),
 \qquad\text{provided }t+h\ne1.
\]
Thus PA194N keeps the consequent in the current formula and offset
registers, preserves the stack on both MP branches, and retains the
original two-formula halt guard. It is neither the original protocol nor
the cached protocol with controls renamed.

For example, take $(t,w,d,h)=(2,5,0,72)$ and parameter $x=3$.
Since $I(2,3)=72$ and $[0;5]=20$, the original controller returns
$(3,20,3,0)$, PA194N returns $(3,5,3,0)$, and the cached controller
returns $(0,5,0,3)$. This is a direct operational distinction, not merely
a difference between table hashes.

\subsection{One literal row and nineteen primitive sites}

In the primary PA194 native table make exactly this replacement:
\[
 \delta(164,1)=(1,R,98)
       \quad\longmapsto\quad (1,R,114).
\]
The old target implements \file{node55 = (node49,node54)}. The new target
starts directly at \file{node54}. Here \file{node49} is the Cantor-stack
push and \file{node54} is the preserving copy of $D$ into $T$ through
scratch register $X$. After successful comparison, $T=H=X=0$ and $D=b$.
The skipped push would therefore have pushed zero. The retained loop
returns $T=b,D=b,X=0$, giving the stated interface.

The retained copy loop has an elementary all-natural proof: after $j\le b$
iterations of its first phase, $(T,D,X)=(j,b-j,j)$; after $j\le b$
iterations of its restoring phase, $(T,D,X)=(b,j,b-j)$. It exits with
$(b,b,0)$ after exactly $7b+2$ primitive source steps.

The source splice removes old primitive ordinals $[3634,3653)$, a block of
19 sites. Its only external entry is the edge $3628\to3634$. The builder
checks every surviving instruction and both targets of each decrement
under the explicit ordinal relocation. The seven-instruction remaining
copy loop is written out in the construction receipt. This binds the
mathematical command change to the actual source, instead of inferring
it from a sampled execution.

\subsection{A simpler exact reset}

A useful consequence of the changed stack behavior is the stronger identity
\[
 \MP(0)(t,w,0,h)=(0,w,0,0)\qquad(t,w,h\in\mathbb N).
\]
Here the output parameter $b=d+0$ is zero. Both the successful and the
failed comparison therefore return exactly $(0,w,0,0)$, regardless of
whether $I(t,0)=h$. Hence one command implements \emph{Discard}, even with
an arbitrary incoming cache. No disjointness assumption about the
truth and implication codes is needed.

Omitting a successful-MP zero push was explored previously at larger state
counts, including in the PA and ZFC material. No priority is claimed for
that general idea. The new result is its complete 194-state realization,
its ten-entry separate-head command compiler and exact physical binding,
and its verified novelty against the current supplied frontier.

\section{A complete ten-entry command compiler}
\label{sec:np-controller}

\subsection{Exact numerical operations and command contracts}

Let
\[
 p(a,b)=\frac{(a+b)(a+b+1)}2+b,\qquad [a;w]=p(a,w),\qquad
 I(a,b)=4p(a,b),\qquad U(i,a)=4p(a,i)+2.
\]
Iterated brackets put their leftmost entry at the top. The Cantor
projections are total and satisfy
\[
 \operatorname{hd}[a;w]=a,\quad \operatorname{tl}[a;w]=w,
 \quad[0;0]=0,\quad\operatorname{hd}(0)=\operatorname{tl}(0)=0.
\]
A numerical stack is therefore not being treated as a free list.
A command store is $(t,w,d,h)$: current formula, stack, offset, and cached
formula. Code zero is truth and code one is the target contradiction.
The new MP command, with $y=d+x$, is
\begin{equation}\label{eq:np-mp}
 \mathsf{MP}_x(t,w,d,h)=
 \begin{cases}
 (y,w,y,0),& I(t,y)=h,\\
 (0,w,y,0),& I(t,y)\ne h.
 \end{cases}
\end{equation}
Shift is retained literally:
\begin{equation}\label{eq:np-shift}
 S(t,w,d,h)=
 \begin{cases}
 \mathsf{halt},&t+h=1,\\
 (h,\operatorname{tl}(w),d,\operatorname{hd}(w)),&t+h\ne1.
 \end{cases}
\end{equation}
For $D=d+x$, $H=h+y$, the unchanged AX command has output
\begin{equation}\label{eq:np-ax}
 \mathsf{AX}_{A,x,y}(t,w,d,h)
   =\bigl(g(H,D,A),[q_1,\ldots,q_9,t;w],0,0\bigr).
\end{equation}
Here $q_1=\Gamma(A)=I(A,U(A,A))$, $q_2,\ldots,q_7$ are the
six retained PA196 heads in reverse emission order,
\[
 q_8=\rho(D,A,H,d)=I(I(I(H,d),D),I(I(D,H),I(A,H))),
 \qquad q_9=U(d,h).
\]
The gate $g$ is the original separate-head guarded separation gate.
The free-constructor computation in
\texttt{check\_emitter\_and\_words.py} binds all these formula trees to
the delivered grammar. All $q_1,\ldots,q_8,g$ are divisible by four,
and $q_9\equiv2\pmod4$, with $q_9\ge2$. In particular none is code one.
This is a statement about constructor values, requiring no consistency
assumption and holding for arbitrary incoming natural stores.

The count in (\ref{eq:np-ax}) is nine generated stack entries plus the
saved old top. The two identities needed below are
\begin{align}
 \mathsf{AX}_{A,x,y};S^9(t,w,d,h)
   &=(q_8,[t;w],0,U(d,h)),\label{eq:np-nine}\\
 \mathsf{AX}_{A,x,y};S^{10}(t,w,d,h)
   &=(U(d,h),w,0,t).\label{eq:np-ten}
\end{align}
All displayed Shifts are nonhalting: their tests involve two generated
non-one codes, zero and a generated code, or the universal code of value
at least two. This also establishes exact stack effects when zero pairs
collapse.

\subsection{Uniform stack operations and a one-command discard}

Write $A_0=\mathsf{AX}_{0,0,0}$ and define
\[
 \mathsf{Push}=A_0;S^9;\mathsf{MP}_0,\qquad
 \mathsf{Pop}=A_0;S^{11},\qquad
 \mathsf{Discard}=\mathsf{MP}_0,\qquad
 \mathsf{Offset}_i=\mathsf{MP}_i.
\]
Their exact contracts are
\begin{align}
 \mathsf{Push}(t,w,d,h)&=(0,[t;w],0,0),\label{eq:np-push}\\
 \mathsf{Pop}(t,[a;w],d,h)&=(t,w,0,a),\label{eq:np-pop}\\
 \mathsf{Discard}(t,w,0,h)&=(0,w,0,0),\label{eq:np-discard}\\
 \mathsf{Offset}_i(0,w,0,0)&=(0,w,i,0).\label{eq:np-offset}
\end{align}
For Push, (\ref{eq:np-nine}) makes the final MP compare an implication
with a universal code, so it fails with offset zero. For Pop, the first
ten Shifts give (\ref{eq:np-ten}). The eleventh Shift tests
$U(d,h)+t\ge2$, hence it does not halt even when $t=1$.
It pops exactly the displayed item. Thus Pop is total on all natural
stores. In particular $\mathsf{Pop}(t,0,d,h)=(t,0,0,0)$ and
$\mathsf{Push}(0,0,d,h)=(0,0,0,0)$.

For Discard, the two MP outcomes in (\ref{eq:np-mp}) coincide when
$d=x=0$, for arbitrary $h$: both are $(0,w,0,0)$. This gives a single-command discard,
where the old dummy-push controller needed a longer word. The old cache
need not be zero for this identity. For Offset,
if $i>0$ then $I(0,i)>0$, so the comparison fails; if $i=0$ it succeeds
and returns $(0,w,0,0)$. The same one-command definition covers both cases.
These simplifications affect the proof-word compiler; they do not claim
additional physical state savings.

\subsection{Axiom selectors and generalization}

In the following contracts the arbitrary incoming store is $(t,w,0,0)$.
For $1\le j\le7$, the word
\[
 \mathsf{Discard};\mathsf{AX}_{A,D,H};S^{j+1};\mathsf{Pop}^{9-j}
\]
returns $(q_j,w,0,0)$ with the parameters evaluated at the desired
$A,D,H$. After the initial Shifts, the remaining $9-j$ stack entries
are exactly the unselected generated tail and the saved zero. The Pops
preserve the selected top and consume those entries; the last pop puts
that saved zero into the cache. For the dominated head, if $u\le a$,
use
\[
 \mathsf{Discard};\mathsf{Offset}_u;
 \mathsf{AX}_{b,a-u,c};S^9;\mathsf{Pop}.
\]
It returns $(\rho(a,b,c,u),w,0,0)$, with the numerical dominance condition
coming from $a=u+(a-u)$. A permitted gate is obtained by
\[
 \mathsf{Discard};\mathsf{AX}_{A,D,H};\mathsf{Pop}^{10},
\]
which returns $(g(H,D,A),w,0,0)$. Truth uses Discard alone.

The word
\[
 \mathsf{Gen}_i=\mathsf{Push};\mathsf{Offset}_i;S;
                \mathsf{AX}_{i,0,0};S^{10}
\]
returns $(U(i,a),w,0,0)$ from $(a,w,0,0)$. Before its AX the store is
$(0,w,i,a)$; its last ten Shifts use (\ref{eq:np-ten}) with saved top zero.
Its first Shift tests $0+0$, and the later ones are among the safe AX
shifts. Hence this generalization word is total even for $a=1$.

\subsection{Modus ponens with its spare zero supplied by the compiler}

Suppose recursively compiled proofs return $a$ and $I(a,b)$ with the
stack framed and both caches cleared, or halt earlier. The MP compiler is
\begin{equation}\label{eq:np-compile}
 \mathsf{Discard};\mathsf{Push};C(I(a,b));\mathsf{Push};
 C(a);\mathsf{Push};S;S;\mathsf{MP}_b;\mathsf{Pop}.
\end{equation}
The first two commands reserve a zero below the premise computations.
If the recursive calls have returned, the store before the two Shifts is
\[
             (0,[a,I(a,b),0;w],0,0).
\]
The first Shift returns $(0,[I(a,b),0;w],0,a)$.
The second halts if $a=1$; otherwise it returns
$(a,[0;w],0,I(a,b))$. MP succeeds and returns
\[
                         (b,[0;w],b,0).
\]
By the total Pop contract, the final word returns $(b,w,0,0)$, including
$b=1$. The stack zero is now consumed once, at the compiler level.
At $a=b=w=0$ the same proof uses the literal Cantor projection identities;
no extra distinction between a zero push and an empty stack is introduced.

Induction on a finite derivation $p$ in the inherited separate-head
calculus therefore proves:
\begin{quote}
From every $(t,w,0,0)$, $C(p)$ either halts or returns
$(\operatorname{concl}(p),w,0,0)$.
\end{quote}
The axiom and generalization cases use the words above. The MP case is
(\ref{eq:np-compile}); extending a word after an early halt still halts.
In particular a refutation compiler ends in top one, if it has not already
halted. Its final appended Shift halts because $1+0=1$.
For arbitrary incoming $t,w,d,h$, first apply $A_0$ to clear $d,h$.
Thus the concrete word
\[
                         A_0;C(p);S
\]
halts from \emph{every} incoming store whenever $p$ derives code one.
The compiler is a primitive-recursive traversal of the finite derivation,
and PA proves its correctness by induction on that coded derivation.
It does not assume a reset between the scheduler's unsuccessful trials.

The unchanged scheduler enumerates finite command words with the usual
Cantor continuation coding. A cofinal family of halting fresh candidates
is obtained by prefixing the universal word by
$\mathsf{AX}_{k,0,0}$. Let $E(\epsilon)=0$ and let a positive continuation be encoded by
$E(c::w)=1+\operatorname{encodeCommand}(c,E(w))$. In particular
$\operatorname{encodeCommand}(\mathsf{AX}_{A,x,y},z)
=p(A+2,p(y,p(x,z)))$. Put $z=E(A_0;C(p);S)$, using this positive
continuation convention rather than a fresh-candidate code. The fresh
candidate codes for the prefixed words are exactly
$p(k+2,p(0,p(0,z)))$, which are strictly increasing and
unbounded in $k$. The fresh counter begins at two, so omitting smaller
candidates loses none of this cofinal family. Every decoded continuation
strictly decreases, and the unchanged source macro loops terminate.
The resulting finite-word and finite-trace construction is formalizable
in PA by primitive recursion and induction.

\subsection{Soundness and the PA calibration}

Annotate the current top $t$, cached head $h$, and every live canonical
Cantor-stack entry by a derivation in the separate-head calculus, or by a
PA proof of its arithmetic translation. For $w>0$, its decoded tail is
strictly smaller than $w$; thus the canonical stack and its annotation
have a primitive-recursive bounded construction. A zero stack supplies
truth when popped, and pushing zero onto zero preserves the empty
canonical annotation.

AX adds the unchanged valid axiom instances, a generalization of the old
cached theorem, and the saved old top. On a successful MP, the annotations
of $t$ and $I(t,y)=h$ produce an annotation for $y$. The stack is retained,
and the cache receives the truth proof. Failure outputs truth and retains
the stack. Shift transports the annotations. If its test succeeds, then
$t+h=1$ implies that one of the two annotated formulas is code one.
The inherited syntactic arithmetic translation turns that annotation into
a PA refutation. No annotation is imposed on the arithmetic offset $d$.

Consequently PA proves both finite syntactic reductions:
\[
 \text{source halt}\ \longrightarrow\ \text{PA refutation},
 \qquad
 \text{PA refutation}\ \longrightarrow\ \text{source halt}.
\]
For the second reduction, compose the supplied primitive-recursive
PA-to-separated-head proof translation with the all-store compiler above.
For the first, compose the proof annotation with the supplied
separated-head-to-PA translation. This uses the existing logical calibration,
including its treatment of the code-zero alias and coincident variable
labels. Nothing in the no-push modification changes that calculus or
requires a new logical axiom, a reflection principle, or consistency as
an additional assumption.



\section{The inherited logical calibration and PA internalization}
\label{sec:logic}

The logical source is the same separate-diagonal-head numerical calculus
as the principal PA194. Its constructors are
\[
 I(A,B)=4P(A,B),\quad M(a,b)=4P(a,b)+1,
 \quad U(a,B)=4P(B,a)+2,\quad E(a,b)=4P(a,b)+3,
\]
with negation $N(B)=I(B,1)$. It retains the six quantified/equality heads,
the numerically dominated propositional schema, theorem generalization,
and modus ponens. The separate diagonal head is
$\Gamma(B)=B\to U(B,B)$; the guarded pointed-separation gate is unchanged.
AX saves ten entries, including the incoming top; the compiler in the
preceding section uses this exact count.

The supplied calibration gives primitive-recursive transformations of
finite refutations in both directions, with correctness provable in PA.
It is a syntactic theorem about proof codes. The soundness direction
interprets membership by finite binary membership, equality by arithmetic
equality, and quantification by first-order quantification. For each fixed
arithmetic formula $\theta$, PA proves finite binary comprehension by
induction on a finite bound. Substituting the syntax of $\theta$ into that
proof template is primitive recursive; no uniform arithmetic truth
predicate or PA reflection principle is asserted.

The converse derives the interfaces used by the original PA refutation
compiler: both membership rewrites, neighborhood equality, and equality
witnesses. Numerical code descent supplies the freshness of label $B$
in its own decoded formula, so the diagonal head replaces the older
stronger equality-vacuity interface. Bundled and separated diagonal gates
have primitive-recursive proof translations. These are the logical
reductions used here, not a substitution of iPRA consistency for PA
consistency.

The supplied argument explicitly treats $I(0,0)=0$, $[0;0]=0$, and
coincident variable labels in pointed separation. In particular, when
the two bound labels coincide, the comprehension matrix is adjusted to
use the singleton condition; treating the original free parameter as
unchanged under that binder would be incorrect. The relevant full supplied
calibration fragments are retained under \file{references/logical/},
with their attribution and exact source snapshot.

Let $R_N$ be the modified finite register program and let $H_R$ express a
finite source halt from the verified clean seed. Command soundness labels
every live formula and canonical stack cell with a finite PA proof of its
translation. The new successful MP removes a push of truth and leaves
the previously annotated stack unchanged. The failed branch creates truth;
Shift, AX, and scheduling retain their old annotation rules. If the sum
of the two visible formula codes is one, one of the annotations is a
proof of the contradiction code. This yields a PA refutation.

Conversely, the inherited PA-to-native proof translation followed by the
new strong command compiler produces a finite halting command word.
Its arbitrary-store contract and cofinal word encoding handle the
scheduler's carried state and any finite initialization prefix. All these
are primitive-recursive finite proof/trace constructions whose correctness
is proved by ordinary induction in PA. Therefore
\[
 \PA\vdash H_R\leftrightarrow
       \exists p\,\Prf_{\PA}(p,\ulcorner0=1\urcorner).
\]
The next section transports this statement to the literal 194-state table.

\section{Physical construction and finite-trace preservation}
\label{sec:physical}

The numerical axiom source, the physical register arrangement, and the
seven simultaneous AX folds are unchanged. The state accounting is
\[
 \underbrace{134}_{\text{fixed controls}}+
 \underbrace{67}_{\text{raw AX readers}}-
 \underbrace{7}_{\text{compatible folds}}=194.
\]
The raw table has 201 controls. The seven folds are reconstructed, and
all 28 context/read-bit contracts are replayed in the final simultaneous
quotient. The check includes incoming transitions establishing the history
bits; the altered MP row does not target a fold member. Full primitive
source simulation checks the resulting table independently of that local
quotient argument.

The logical register order is $(N,T,W,A,D,H,X,Y,L)$, with physical map
$(6,2,0,4,3,7,8,5,1)$. A block $0\,1^{v+1}$ encodes natural value $v$.
The source frame contains the left marker, the source-path bits, an affine
zero gap, the delimiter, the nine blocks, and an arbitrary finite nonempty
list of spare blocks. The header distance is $48+n$, for arbitrary
$n\ge0$. The increment operation's append-one spare effect is included.

Exact blank execution reaches source site 25, native state 84, with head
$-2$ and logical registers $(2,0,0,0,0,0,0,0,0)$, after 58,154 transitions.
The tape check compares its entire finite support, every source bit and
marker, all nine register blocks, all 26 one-valued spare blocks, and
zero exterior. The source gap invariant starts at this exact frame.

\begin{lemma}[Positive primitive simulation]
Each nonhalting source instruction of $R_N$ has a positive, terminating
literal-tape realization from every permitted source frame. It reaches
exactly the prescribed successor frame with the specified register and
spare-tail updates. Both branches of every decrement are covered. A source
halt reaches the physical terminal control.
\end{lemma}
\begin{proof}
The delivered finite instruction map has 3,654 sites and 4,755 guarded
outcomes. The construction checker verifies every one using the fixed
register-engine contracts. The independent checker derives four generic
list-sweep lemmas from the actual transition rows and checks the same
outcomes using a different literal-tape representation and independently
resolved control flow. Its proof language consists of literal steps,
constant-symbol scans, and generic finite-list sweeps. Induction on unary
block length and spare-list length justifies those rules. Instantiating
the finite derivations proves the contracts for arbitrary frame parameters.
Every segment has positive duration; the smallest certified affine
constant is three. These elementary induction arguments and finite
certificate instances are formalizable in PA.
\end{proof}

\begin{theorem}[PA-provable source/tape halting equivalence]
\[
             \PA\vdash H_N\leftrightarrow H_R.
\]
\end{theorem}
\begin{proof}
A halting source trace gives a halting tape trace by concatenating its
primitive simulations after the exact startup. Conversely, if the tape
first halts after $T$ steps, iterate the total source step function for
$T+1$ instructions from the clean seed, making a source halt absorbing.
If no source halt occurs, positive primitive simulations concatenate to
a nonhalting tape prefix of length greater than $T$, contradicting
machine determinism. All iterations and finite trace concatenations are
primitive recursive. PA proves their correctness by induction; the
argument does not assume the actual run is nonhalting.
\end{proof}
Composing with the preceding logical/source equivalence proves the theorem
in Section~\ref{sec:result}.

\section{Verification results and exact scope}
\label{sec:verification}

\begin{center}
\begin{tabular}{@{}lr@{}}
\toprule
Check & Observed result\\
\midrule
Working controls / literal rows & 194 / 388\\
Raw AX readers / compatible folds & 67 / 7\\
Primitive source sites / outcomes & 3,654 / 4,755\\
Construction-checker symbolic steps & 1,473,292\\
Independent-checker symbolic steps & 1,073,769\\
Independent maximum per outcome & 397\\
Independent minimum duration constant & 3\\
Generic loop lemmas rederived from rows & 4\\
Exact complete-tape blank startup & 58,154 steps\\
Previously represented 194-state graph classes & 8\\
Previously represented 194-state byte forms & 15\\
\bottomrule
\end{tabular}
\end{center}

Symbolic-step totals measure verifier work, not a uniform literal running
 time for arbitrary natural register inputs. Startup is a finite exact
initialization check; it is not empirical evidence for infinite nonhalting.
The universal contracts cover nine independent natural register values,
an arbitrary affine gap and an arbitrary nonempty finite spare list.
The two checkers are independently implemented at the physical level;
they are not two independently developed logical PA calibrations.

The corpus novelty checker compares exact canonical transition encodings,
not only differing file hashes. The new root-free graph hash is
\begin{quote}\small
\nolinkurl{e74434eb401367ab63d712bbde26a0edcec17b2ac536ca4d474e743d9969b036}.
\end{quote}
It differs from every recognized existing 194-state graph even after
arbitrary state renaming, changing a graph-reaching root, tape reflection,
and simultaneous read/write symbol complementation. The last symmetry
is excluded only as cosmetic novelty; it would alter the specified blank
symbol if used as a machine transformation.

The mathematical statement is an ordinary PA-provability theorem built
from the explicit proof transformations and physical induction arguments.
The field \file{kernel_checked: false} records that no new end-to-end
Lean/Rocq theorem was compiled. The repository's open machine-specific
iPRA obligations remain separate. No Python run alone proves the inherited
logical calibration.

\section{Comparison with the complete supplied corpus}
\label{sec:corpus-review}

The comparison is pinned to ProveIt commit
\nolinkurl{17f92105e2bc4317feae6c219cb2143437b482ff}.
At this revision the PA research directory contains the extracted contents
of 36 delivered packets. The incoming directory contains 26 ZIP files:
21 are retained copies of those packets, and five are later cached-controller
packets. Thus there are 41 top-level research packets to compare, rather
than 62 independent constructions. The numerous
\nolinkurl{pa194-*} and \nolinkurl{pa195-*} proof directories develop
different proof layers for the same designated machines; their names do
not signify additional machine designs.

The comparison has two levels. First, a fresh read-only census scans every
\nolinkurl{.tm} and \nolinkurl{.txt} file under the entire PA directory,
every incoming ZIP, and every ZIP recursively encountered in those sources.
Archive payloads are deduplicated by their complete SHA-256. Second, the
reports and source descriptions classify the recognized tables as principal
results, supplementary implementations, search seeds, or rejected
experiments. A small state count or a filename beginning with ``candidate''
is never sufficient for a positive classification.

The census visits 44 distinct ZIP payloads through 71 archive references.
It recognizes 4,417 table occurrences, 620 exact byte identities, and 446
action-preserving, root-free control-graph classes. Two disconnected
synthetic controls have no root reaching their whole graph and are left
unclassified. The numeric tables contain eight 194-state graph classes,
sixteen 195-state classes, and ninety-nine 196-state classes. These are
corpus counts, including diagnostics and repeats of ancestors; they are
not counts of established consistency machines. No recognized complete
table has between 10 and 193 working states. This observation is not a
state-count lower bound.

\subsection{The eight existing 194-state graphs}

\begin{longtable}{@{}p{0.12\linewidth}p{0.20\linewidth}p{0.58\linewidth}@{}}
\toprule Family & Controls / source sites & Change and status\\\midrule\endhead
PA194-67 & 134 + 67 - 7; 3673 sites & Separate diagonal head; 67 AX readers and seven compatible history folds. principal; ordinary PA proof and two symbolic physical checkers\\[4pt]
PA194-AD & 134 + 67 - 7; 3673 sites & Exchange the physical A and D register positions; three changed native rows. principal companion; ordinary transfer and two symbolic physical checkers\\[4pt]
PA194-68 & 134 + 68 - 8; 3673 sites & Different association of the same separate-head emitter, with eight folds. principal; ordinary PA proof with independently reconstructed source/table\\[4pt]
PA194C & 134 + 67 - 7; 3649 sites & Cached Shift/MP controller on PA194-67; six changed native rows. All five later cached packets duplicate this exact dense table. principal; ordinary inherited PA calibration and two symbolic physical checkers\\[4pt]
PA194CS & 135 + 67 - 7 - 1; 3646 sites & Destructive equality operand symmetry shortens Shift by three sites. Extra controller reader 330 is folded with 41; this is not an extra AX fold. principal; source block equivalence and two symbolic physical checkers\\[4pt]
PA194altC & 134 + 68 - 8; 3649 sites & Cached controller composed with the alternate PA194-68 reader allocation. search seed; no separate complete consistency endpoint delivered\\[4pt]
PA194CZ & 134 + 67 - 7; 3649 sites & Replace source-site-26 jump by a zero test of register six; native row (27,1) changes from 0 L 40 to 1 R 21. supplementary seed; primary contracts and positivity-mask invariant, no new independent physical replay\\[4pt]
Rejected scanner194 & rejected & Optimistic shortened-scanner/reader quotient fails the required source simulation. explicitly rejected frame-contract experiment; not a Con(PA) result\\[4pt]
\bottomrule
\end{longtable}


The five principal designs in the table are PA194-67, PA194-AD, PA194-68,
PA194C, and PA194CS. The remaining three are deliberately retained in the
novelty comparison: PA194altC and PA194CZ already occur as supplementary
search seeds, while the shortened-scanner table is explicitly rejected.
All eight classes remain distinct even if one allows tape reflection,
simultaneous read/write symbol complementation, and arbitrary choice of
a graph-reaching root. Symbol complementation is used only as a stronger
test against cosmetic novelty; it does not preserve the specified all-zero
blank input.

Several archive names suggest different cached machines, but their
principal literal result is the same table. Its SHA-256 is
\begin{quote}\small
\nolinkurl{d759313c03b515293fe06b12769fa1771d3829a2bad89febaa9b4391990f3a4b}.
\end{quote}
In particular, recomposing the primary seven-fold decoder with the
Combined controller would reproduce an existing result. The symmetry
of the destructive equality test is also already used by PA194CS.
Likewise, merely installing the Combined controller on the alternate
eight-fold decoder, or adding the zero-register jump at source site 26,
would reproduce supplied search seeds. A proposed fallback must be
compared with these supplementary tables as well as the principal ones.

\subsection{The positive 195-state families and their limitations}

\begin{center}\small
\begin{tabular}{@{}lrrrl@{}}
\toprule
Family & Sites & Outcomes & Startup & Physical organization\\
\midrule
Original PA195 & 3,700 & 4,815 & 58,154 & $134+68-7$\\
195S & 3,673 & 4,780 & 58,154 & $134+67-6$\\
195C & 3,649 & 4,749 & 57,922 & 195S with Combined commands\\
M195S8 & 3,700 & 4,815 & 59,164 & Extra scanner control and eight folds\\
PA195DC & 3,676 & 4,784 & 57,922 & Original diagonal emitter, Combined commands\\
\bottomrule
\end{tabular}
\end{center}

The table identifies the five designated positive manuscript families.
Original PA195 packages the diagonal-vacuity formula with the guarded
gate. The separate-head emitter saves one macro occurrence and 27
primitive sites; it underlies 195S, 195C, and the principal 194-state
families. Combined Shift and MP remove a further 24 primitive sites,
but do not automatically save a reader state. M195S8 uses an additional
scanner control to make a different eight-fold allocation valid.
PA195DC applies the Combined controller to the original bundled-diagonal
source rather than the separate-head source.

The sixteen recognized 195-state graph classes include these five, six
explicitly rejected or diagnostic quotients, a secondary separate195
layout, a crossed allocation, the 195-state parent of PA194CS, and two
secondary implementations in the PA195DC search material. The latter
include a negative-diagonal variant whose author explicitly keeps it
outside the principal consistency theorem. They must not be silently
promoted to independently proved consistency constructions. The full
class inventory records every byte identity and source locator.

\subsection{What the larger designs contribute}

James Chang's PA209 and the Joint/Compatibility U209 ancestor are different
209-state constructions, with respectively 3,948 and 3,909 primitive
source sites. The U208 packets identify two old decoder roles by a
finite-history argument. Entry-context, three-cell, and history
presentations repeatedly describe this same kind of one-state saving;
later helper layouts produce further 208-state graphs. Their equal
state counts do not identify their source programs or transition graphs.

The published PA196 has 3,812 source sites and 4,959 primitive outcomes.
Its allocation comprises 36 tape/register-engine controls, 13 primitive
routers, and 147 readers. The later 196-state proposals change head
ordering, association, physical A/D positions, closing decrements, or
zero-register jump routing. The direct repository catalogue selects
22 rooted 196-state graph classes and explicitly distinguishes two
experimental MP variants and two rejected layouts. The latter both
halt on blank tape after 43,563 transitions; their finite halting traces
are Lean checked in the repository. They are controls, not consistency
machines.

The mathematical source changes show why source length cannot be used
as a proxy for the number of physical states. Diagonal vacuity reduces
the emitter to 3,700 sites, but its first implementation uses 199 states;
later allocations reach 198, 196, and 195. Separating the diagonal head
gives 3,673 sites and eventually reaches 194. An integrated witness gate
uses 3,745 sites and 4,872 primitive outcomes, yet its verified physical
allocations use 206 and then 203 states. Distributed-gate and
short-negation alternatives likewise lose advantageous reader sharing.
The NEG18 separate-emitter source has only 3,441 sites and 4,468
outcomes, but the documented allocation is $134+68-5=197$.
A repaired 157-node representation of the unchanged PA196 source has
a 207-state realization: a smaller abstract DAG can require more
routing and fewer useful contextual identifications.

Changed command protocols form a separate axis. H-only Shift alone has
3,811 sites; Combined196 has 3,788 sites, 4,928 outcomes, and a complete
57,922-step startup. Its successful MP returns the consequent in H,
clears T and D, and does not push a dummy stack value. Shift tests H.
These semantics differ from the original controller, so the command
compiler and arbitrary-incoming-store refutation argument must change
with them. An MP-only bounded register experiment is not the same result
as a complete Combined implementation with all universal source outcomes.

Other packets contain source-only grammar proposals, failed scanner
shortcuts, pruned sources without a calibrated blank seed, and
intermediate parity/control components. Their contribution is evidence
about particular search spaces or reusable proof obligations. A failed
universal frame contract rejects the proposed simulation argument; it
does not by itself prove that the failing frame occurs on the blank run.
Conversely, a long nonhalting prefix does not repair that argument.

\subsection{Proof status and inherited assumptions}

The ordinary PA calibration in the recent reports is stronger than
experimental evidence but distinct from a freshly elaborated proof-assistant
endpoint. It supplies syntactic translations of PA and numerical-calculus
refutations, an explicit command compiler, exhaustive command enumeration,
and positive-duration physical simulations. The reports explicitly address
the code-zero alias $I(0,0)=0$, the Cantor-stack alias $[0;0]=0$, early
halts during proof compilation, an arbitrary carried store between trials,
and the coincident-variable case in the finite bit-comprehension argument.
These points matter to the claimed PA-provable biconditional.

The repository's current README nevertheless states that neither
machine-specific reduction has yet been established in iPRA for the
primary PA194, original PA195, or PA196 candidate. Its ordinary finite-run
Lean theorems and conditional arithmetic endpoints do not close those
obligations. This is consistent with reusing an explicit ordinary
mathematical PA calibration; it must not be described as a completed
Lean/Rocq or iPRA result. The present comparison audits the arguments and
their interfaces but does not independently rebuild every historical
proof-assistant development.

\subsection{Coverage limits and reproducibility}

The literal recognizer requires a UTF-8 text file of five-column rows
\texttt{state read write direction next}, numeric state labels, a binary
alphabet, both rows for every working state, and exactly one undefined
halt target. Blank lines and whole-line comments are allowed. Files using
other syntaxes, multiple undefined halts, code-generated tables not
materialized as text, and tables embedded solely in JSON or Lean are
outside this mechanical recognizer. They are considered through the
reports and source descriptions rather than counted as recognized literal
tables. The archive ledger gives all 41 top-level packet contributions.

Graph classification forgets the designated start and compares each
root that reaches the full working graph by deterministic breadth-first
numbering with symbol-labelled edges and exact write/direction actions.
Taking the least such encoding yields a complete invariant for this
stated class of root-free graphs. This is stronger than comparing native
labels or file hashes, but makes no assertion about semantic equivalence
or correctness. Exact machine bytes, graph hashes, all occurrence paths,
archive identities, the parser, and the read-only reproduction command
are included in the comparison artifacts.

\subsection{Prior work on expanded macro searches}

The supplied search material already contains a partial expansion of
the parameter-copy macros. In
\nolinkurl{pa194_conpa_2026-11-13_2/search/exploratory/expanded_input.txt},
fourteen terminals replace the eleven-terminal basis. Native copy roots
12, 98, and 166 are opened into push root 19 and tail roots 20, 114,
and 179. The old diagonal macro word grows from 191 to 252 terminals.
The associated \nolinkurl{search_expanded.cpp} hashes newly composed
pairs against all fixed pure-reader pairs, so exact fixed-pair reuse
is already present. The recorded expanded and general runs cover
3,363,000 and 1,023,000 iteration slots, respectively, and retain 195
states as their best count. Earlier free-macro-aware grammar searches
also give source-only upper bounds.

Opening deeper interiors of the cons and push macros, using the newer
separate-head or cached source, and performing a physically checked
global optimization can extend those searches. It would be inaccurate
to claim that exposing the three copy tails or recognizing fixed pure
reader pairs is itself a new idea. Nor do those negative runs prove
that a deeper or differently initialized search cannot save a state.

\section{Search below 194 and its limits}
\label{sec:search}

The positive deliverable is PA194N. No smaller machine was certified.
The new search effort addressed two concrete limitations of the supplied
searches: treating large fixed subroutines as indivisible leaves, and
requiring every incoming edge of an eliminated state to share one replacement.
It also re-examined an outer-controller rotation motivated by a possible
shorter return scanner.

\subsection{Deeper macro expansion and broader fold allocation}

The first decoder experiment opens one pure-reader level of every eligible
axiom macro, including the Cantor push and constructor routines as well as
the three parameter copies. A deeper experiment opens a second level of
the push, constructor and guard routines while retaining the copy-tail
interfaces. These extend the earlier three-copy-tail experiment identified
in the corpus review. The separate-head source and its new no-push controller
supply different starting points from that earlier diagonal source.

The search interns new binary reader pairs against the fixed controller's
existing pure-reader pairs. It uses both whole-shared-subtree and individual
occurrence rotations. In broader runs, a candidate fold may identify an AX
reader with a fixed control, provided all incoming rows establish the
required preceding bit. The local simulation is checked in the simultaneously
patched table. One heuristic variant prunes partial fold selections that
already fail their local check; that is a search restriction, not a completeness
theorem about possible quotients.

\begin{center}\small
\begin{tabular}{@{}lrrr@{}}
\toprule
Search & Seconds & Proposal slots & Best score\\
\midrule
One-level expansion; AX folds & 600 & 6,548,653 & 194\\
One-level expansion; fixed/AX folds & 600 & 3,584,338 & 194\\
Partial-selection pruning & 600 & 3,487,411 & 194\\
No-push controller; one level & 700 & 5,228,572 & 194\\
No-push controller; deeper interiors & 700 & 6,366,489 & 194\\
\midrule
Total & 3,200 & 25,215,463 & 194\\
\bottomrule
\end{tabular}
\end{center}


An iteration is a mutation-proposal slot; it need not denote a distinct
candidate, and some slots produce no mutation. The reported best value is
only a screened allocation score, equal to 134 plus reader count minus
compatible folds. No value below 194 was reached. A grammar achieving a
promising score would still have to pass exact source-graph binding,
complete physical simulation, and calibrated blank initialization. These
heuristic searches therefore establish neither optimal grammar size nor
a state-count lower bound. The accepted PA194N does not depend on their
termination or on the optimality of their result.

The frozen inputs, source maps, search sources, random seeds, complete logs
and terminal iteration totals are included. Time-limited runs can have
different iteration counts on different computers; the retained data,
not a promise of deterministic stochastic discovery, defines the evidence.
\subsection{Additional physical contraction searches}
Two further searches were carried out on the existing PA194CS table and
on the new no-push table. The contexts were regenerated from all source
outcomes, symbolic unary-sweep interiors, generic list-block proofs, and
the exact blank initialization. Every window has radius five, and every
local replacement has a 300-transition bound.
\[
\begin{array}{lrrr}
\text{input}&\text{source outcomes}&\text{symbolic views}&\text{windows}\\\hline
\text{PA194CS}&4745&19067&29035\\
\text{PA194N}&4755&19043&29017
\end{array}
\]
The supplied common-action merger yielded no additional state
identification on either set. A new \emph{incoming-edge bypass} search also
yielded none. The latter does not insist on one common survivor for all
incoming edges of a deleted control: it permits each incoming transition
to receive its own new action, provided its bounded continuation reproduces
the old incoming step followed by the deleted control's step. The continuation
must preserve the entire window, avoid the deleted control and all rewritten
rows, and end at the exact old control and head position. Self-looping
removed controls are excluded from this search rule.

These are bounded searches through sufficient rewrite rules. They are not
lower bounds for binary Turing machines and do not exclude larger simultaneous
rewrites or a different register engine. The frozen tables, context files,
search sources and empty candidate outputs with progress logs are supplied
in the supplementary search bundle.

\paragraph{A startup/scanner obstruction.}
Reassociating the outer source tree from
$((G,D),C)$ to $(G,(D,C))$ removes one leading zero from the first
comparison path. Its resolved source graph changes only the final
relative continue target, which would need adjustment from three to four.
However, the obvious four-row physical implementation is invalid: the
shared native control 69 implements the decoder block only with a
preceding one. The rotated entry supplies a preceding zero and selects
its other shared behavior. The candidate fails the universal contract at
source site 11, and its blank startup does not reach the intended clean
frame. This is a rejected experiment, not a 193-state result.

\subsection{What transfers from ZFC and concise iPRA}

The ZFC work supplied concrete precedents for removing the successful-MP
zero push, proving proof compilers over arbitrary retained stacks, and
separating a source change from its physical state price. PA194N uses that
interface discipline and keeps the PA-calibrated numerical calculus.
Its one-command reset is an exact consequence of the altered MP branches;
it is not an additional logical axiom.

The concise-iPRA project provides useful examples of deriving apparently
necessary rule cases from a smaller interface: one-term transport,
restricted-source transport, and combined introduction rules. That suggests
examining special cases and interface redundancy before shortening code.
It does not license replacing the PA consistency search by an iPRA
consistency search. A new smaller axiom basis would require its own
PA-verifiable proof translations and physical compilation. No such
replacement was established here.

The largest remaining opportunities are a jointly changed fixed controller
and AX layout, a new source interface with a smaller physical price, or a
correct global return-scanner redesign. The outer-rotation failure shows
why a syntactically shorter path cannot simply be substituted into a
history-sensitive decoder. These are open directions, not promises of a
193-state construction.


\section{Reproduction and package organization}

Python 3.10 or later and the standard library suffice for the principal
construction, command certificates, and independent physical check.
Assertions must remain enabled. The root \file{README.md} gives exact
commands, and \file{verify_all.py} runs the accepted checks from any
current directory without network access. Search replay alone requires
a C++17 compiler and is not necessary to reconstruct the accepted table.
A standard TeX Live installation builds this report with two passes of
\file{pdflatex} from the \file{report/} directory.

The archive separates accepted construction, independent verification,
mathematical proofs, comparison, inherited references and exploratory
search artifacts. Rejected or unproved search proposals are not used in
the theorem. Exact bytes, snapshot identities and licensing notices are
recorded with the artifacts. No repository modification or publication
is part of this delivery.

\appendix
\section{Complete transition table}
\label{app:table}
An entry $wDq$ writes bit $w$, moves in direction $D\in\{L,R\}$ and
enters $q$. $H$ denotes the separate halt label 194. The start is state
zero on the all-zero tape. The appendix is generated from the literal
\file{pa194N.tm} file.
\begingroup\small\setlength{\tabcolsep}{7pt}
\begin{longtable}{@{}rll@{\hspace{20pt}}rll@{\hspace{20pt}}rll@{}}
\toprule
$q$ & Read 0 & Read 1 & $q$ & Read 0 & Read 1 & $q$ & Read 0 & Read 1\\
\midrule\endhead
0 & \texttt{1R1} & \texttt{1R2} & 65 & \texttt{0R91} & \texttt{1R49} & 130 & \texttt{0R134} & \texttt{1R135}\\
1 & \texttt{1R3} & \texttt{1R1} & 66 & \texttt{0R92} & \texttt{0L75} & 131 & \texttt{0R136} & \texttt{1R137}\\
2 & \texttt{0R0} & \texttt{1R2} & 67 & \texttt{0R86} & \texttt{1R33} & 132 & \texttt{0R138} & \texttt{1R139}\\
3 & \texttt{0R4} & \texttt{0R1} & 68 & \texttt{0R93} & \texttt{1R68} & 133 & \texttt{0R140} & \texttt{1R141}\\
4 & \texttt{1R5} & \texttt{1L6} & 69 & \texttt{0L94} & \texttt{1L33} & 134 & \texttt{0R100} & \texttt{1R67}\\
5 & \texttt{1L7} & \texttt{1L8} & 70 & \texttt{0L84} & \texttt{0R2} & 135 & \texttt{0R142} & \texttt{1R127}\\
6 & \texttt{0L7} & \texttt{1L6} & 71 & \texttt{0R95} & \texttt{1R71} & 136 & \texttt{0R143} & \texttt{1R144}\\
7 & \texttt{0L9} & \texttt{1L6} & 72 & \texttt{0R75} & \texttt{1R96} & 137 & \texttt{0R145} & \texttt{1R146}\\
8 & \texttt{0L10} & \texttt{0L11} & 73 & \texttt{0R59} & \texttt{0R85} & 138 & \texttt{0R134} & \texttt{1R147}\\
9 & \texttt{1R12} & \texttt{0L9} & 74 & \texttt{0R86} & \texttt{1R97} & 139 & \texttt{0R120} & \texttt{1R148}\\
10 & \texttt{0L9} & \texttt{0L10} & 75 & \texttt{0R42} & \texttt{0L13} & 140 & \texttt{0R149} & \texttt{1L131}\\
11 & \texttt{1L8} & \texttt{1L11} & 76 & \texttt{0R87} & \texttt{1R37} & 141 & \texttt{0R150} & \texttt{1R151}\\
12 & \texttt{0R12} & \texttt{0R13} & 77 & \texttt{0R30} & \texttt{1R98} & 142 & \texttt{0R79} & \texttt{1R41}\\
13 & \texttt{0L14} & \texttt{0R13} & 78 & \texttt{0R42} & \texttt{1R99} & 143 & \texttt{0R28} & \texttt{1R152}\\
14 & \texttt{1L15} & \texttt{0L15} & 79 & \texttt{0R100} & \texttt{1R101} & 144 & \texttt{0L153} & \texttt{1L154}\\
15 & \texttt{0L15} & \texttt{1L16} & 80 & \texttt{0R102} & \texttt{1L103} & 145 & \texttt{0L129} & \texttt{1R155}\\
16 & \texttt{0L17} & \texttt{1L16} & 81 & \texttt{0R30} & \texttt{1R80} & 146 & \texttt{0R156} & \texttt{1R157}\\
17 & \texttt{0L18} & \texttt{1L16} & 82 & \texttt{0R97} & \texttt{1R104} & 147 & \texttt{0R149} & \texttt{1R158}\\
18 & \texttt{0L19} & \texttt{1L16} & 83 & \texttt{0R30} & \texttt{1R105} & 148 & \texttt{0R159} & \texttt{1R120}\\
19 & \texttt{0L20} & \texttt{1L16} & 84 & \texttt{1R84} & \texttt{1L15} & 149 & \texttt{0R100} & \texttt{1R160}\\
20 & \texttt{0L21} & \texttt{1L16} & 85 & \texttt{0R106} & \texttt{1R85} & 150 & \texttt{0R161} & \texttt{1R162}\\
21 & \texttt{0L22} & \texttt{1L16} & 86 & \texttt{0R73} & \texttt{1R107} & 151 & \texttt{0R163} & \texttt{0L10}\\
22 & \texttt{0L23} & \texttt{1L16} & 87 & \texttt{0L46} & \texttt{0R93} & 152 & \texttt{0R164} & \texttt{1R165}\\
23 & \texttt{0R23} & \texttt{1R24} & 88 & \texttt{0R108} & \texttt{1R98} & 153 & \texttt{0R166} & \texttt{1R157}\\
24 & \texttt{0R25} & \texttt{1R26} & 89 & \texttt{0R108} & \texttt{1R69} & 154 & \texttt{0R167} & \texttt{1R156}\\
25 & \texttt{0R27} & \texttt{1R28} & 90 & \texttt{0R109} & \texttt{0L14} & 155 & \texttt{0R135} & \texttt{1R130}\\
26 & \texttt{0R29} & \texttt{1R30} & 91 & \texttt{0R110} & \texttt{1L74} & 156 & \texttt{0R147} & \texttt{1R41}\\
27 & \texttt{0R31} & \texttt{1RH} & 92 & \texttt{0R111} & \texttt{1R112} & 157 & \texttt{0R134} & \texttt{1R168}\\
28 & \texttt{0L32} & \texttt{1R33} & 93 & \texttt{0R113} & \texttt{1R93} & 158 & \texttt{0R41} & \texttt{1L169}\\
29 & \texttt{0R34} & \texttt{1R35} & 94 & \texttt{0R50} & \texttt{1R103} & 159 & \texttt{0L170} & \texttt{1R171}\\
30 & \texttt{0L36} & \texttt{1L37} & 95 & \texttt{0R114} & \texttt{1R95} & 160 & \texttt{0R91} & \texttt{1R64}\\
31 & \texttt{0R38} & \texttt{1R39} & 96 & \texttt{0R30} & \texttt{1R49} & 161 & \texttt{0R172} & \texttt{1R170}\\
32 & \texttt{0R40} & \texttt{1R41} & 97 & \texttt{0R48} & \texttt{1R61} & 162 & \texttt{0R173} & \texttt{1R174}\\
33 & \texttt{0R42} & \texttt{1R43} & 98 & \texttt{0R87} & \texttt{1L115} & 163 & \texttt{0R175} & \texttt{1R176}\\
34 & \texttt{0R44} & \texttt{1R45} & 99 & \texttt{0R104} & \texttt{0L14} & 164 & \texttt{0R177} & \texttt{1R178}\\
35 & \texttt{0R46} & \texttt{1R47} & 100 & \texttt{0R36} & \texttt{1R116} & 165 & \texttt{0L179} & \texttt{1L180}\\
36 & \texttt{0R48} & \texttt{1R49} & 101 & \texttt{0R115} & \texttt{1R78} & 166 & \texttt{0R149} & \texttt{1R140}\\
37 & \texttt{0R50} & \texttt{1R51} & 102 & \texttt{0L46} & \texttt{0R113} & 167 & \texttt{0R165} & \texttt{1R181}\\
38 & \texttt{0R52} & \texttt{1R42} & 103 & \texttt{0L45} & \texttt{1R117} & 168 & \texttt{0L130} & \texttt{1R120}\\
39 & \texttt{0R53} & \texttt{1R54} & 104 & \texttt{0L94} & \texttt{1R118} & 169 & \texttt{0R182} & \texttt{1R135}\\
40 & \texttt{0R55} & \texttt{1R56} & 105 & \texttt{0R119} & \texttt{0L14} & 170 & \texttt{0R183} & \texttt{1R145}\\
41 & \texttt{0R57} & \texttt{1R58} & 106 & \texttt{0R71} & \texttt{1R106} & 171 & \texttt{0R158} & \texttt{1R130}\\
42 & \texttt{0R59} & \texttt{0R60} & 107 & \texttt{0R120} & \texttt{1R30} & 172 & \texttt{0R139} & \texttt{1R165}\\
43 & \texttt{0R61} & \texttt{0L14} & 108 & \texttt{0L46} & \texttt{0R106} & 173 & \texttt{0R158} & \texttt{1R184}\\
44 & \texttt{0R62} & \texttt{1R63} & 109 & \texttt{0R70} & \texttt{0R95} & 174 & \texttt{0R185} & \texttt{1R91}\\
45 & \texttt{0R55} & \texttt{1R64} & 110 & \texttt{0R59} & \texttt{0R95} & 175 & \texttt{0R31} & \texttt{1R53}\\
46 & \texttt{0L65} & \texttt{1R66} & 111 & \texttt{0L65} & \texttt{1R121} & 176 & \texttt{0R186} & \texttt{1R175}\\
47 & \texttt{0R36} & \texttt{1R67} & 112 & \texttt{0R122} & \texttt{1R123} & 177 & \texttt{0R134} & \texttt{1R142}\\
48 & \texttt{0R59} & \texttt{0R68} & 113 & \texttt{0R14} & \texttt{1R113} & 178 & \texttt{0R32} & \texttt{1R187}\\
49 & \texttt{0R69} & \texttt{0L14} & 114 & \texttt{0R68} & \texttt{1R114} & 179 & \texttt{0R144} & \texttt{1R188}\\
50 & \texttt{0R70} & \texttt{0R71} & 115 & \texttt{0R124} & \texttt{1L67} & 180 & \texttt{0R189} & \texttt{1R190}\\
51 & \texttt{0R72} & \texttt{0L49} & 116 & \texttt{0R125} & \texttt{1R126} & 181 & \texttt{0R41} & \texttt{1R191}\\
52 & \texttt{0R73} & \texttt{1R54} & 117 & \texttt{0R120} & \texttt{0L14} & 182 & \texttt{0R168} & \texttt{1R41}\\
53 & \texttt{0L74} & \texttt{0L15} & 118 & \texttt{0R94} & \texttt{1L98} & 183 & \texttt{0R192} & \texttt{1R120}\\
54 & \texttt{0R75} & \texttt{0L10} & 119 & \texttt{0R127} & \texttt{1R127} & 184 & \texttt{0R153} & \texttt{1R146}\\
55 & \texttt{0R76} & \texttt{1R77} & 120 & \texttt{0R70} & \texttt{0R68} & 185 & \texttt{0R165} & \texttt{1R180}\\
56 & \texttt{0R78} & \texttt{1R79} & 121 & \texttt{0R128} & \texttt{1R129} & 186 & \texttt{0R96} & \texttt{1R115}\\
57 & \texttt{0R80} & \texttt{1R81} & 122 & \texttt{0R32} & \texttt{1R36} & 187 & \texttt{0R193} & \texttt{1R159}\\
58 & \texttt{0R82} & \texttt{1R83} & 123 & \texttt{0R31} & \texttt{1R101} & 188 & \texttt{0R169} & \texttt{1R154}\\
59 & \texttt{0L84} & \texttt{0R1} & 124 & \texttt{0R59} & \texttt{0R114} & 189 & \texttt{0R158} & \texttt{1R192}\\
60 & \texttt{0R85} & \texttt{1R60} & 125 & \texttt{0R102} & \texttt{1R117} & 190 & \texttt{0R179} & \texttt{1R182}\\
61 & \texttt{0L76} & \texttt{1L86} & 126 & \texttt{0R82} & \texttt{1R81} & 191 & \texttt{0R182} & \texttt{1R41}\\
62 & \texttt{0R87} & \texttt{0L7} & 127 & \texttt{0R120} & \texttt{1R120} & 192 & \texttt{0R149} & \texttt{1R168}\\
63 & \texttt{0R88} & \texttt{1R89} & 128 & \texttt{0R130} & \texttt{1R131} & 193 & \texttt{0R79} & \texttt{1R144}\\
64 & \texttt{0R42} & \texttt{1R90} & 129 & \texttt{0R132} & \texttt{1R133} &  &  & \\
\bottomrule
\end{longtable}
\endgroup


\section{Archive-by-archive ledger}
\begin{longtable}{@{}p{0.40\linewidth}p{0.55\linewidth}@{}}
\toprule Packet & Reported contribution\\\midrule\endhead
\nolinkurl{pa_unary_power_cover} & Unary power covering at PA strength (report of 27 October 2026). The U208 packets carry it as a baseline.\\[4pt]
\nolinkurl{U208_ConPA_Conditional_Context_2026-10-29} & U208 by a conditional tape-context reduction. States a 208-state candidate \file{machine/u208.tm}.\\[4pt]
\nolinkurl{U208_ConPA_Entry_Context_2026-10-29} & U208 by eliminating a state using entry context. States a 208-state machine \file{machine/u208.tm}.\\[4pt]
\nolinkurl{U208_ConPA_History_2026-10-29} & U208 by a finite-history decoder fold. States a 208-state candidate \file{machine/u208.tm}.\\[4pt]
\nolinkurl{U208_ConPA_2026-10-29} & U208 by a three-cell fold. States a 208-state candidate \file{machine/u208.tm}.\\[4pt]
\nolinkurl{u208_research_2026-10-29} & U208 research release. States two 208-state machines, \file{machine/u208.tm} and \file{machine/variant_plain/u208.tm}.\\[4pt]
\nolinkurl{Below_U208_Research_2026-10-29} & Research toward fewer than 208 states. States that no machine with 207 states or fewer was obtained.\\[4pt]
\nolinkurl{u208_comparison_research_2026-10-29} & Comparison of the 208-state proposals and a search below 208. States that the target was not achieved.\\[4pt]
\nolinkurl{ConPA_below208_research_2026-10-29} & Comparative audit and reduction search below 208. States that it does not establish a 207-state machine.\\[4pt]
\nolinkurl{pa_consistency_research_2026-10-29} & Audit and new source reductions. States that no validated machine below 208 states was obtained.\\[4pt]
\nolinkurl{pa208_research_2026-10-29} & Comparison and research below 208 states. States that it does not establish a valid machine below 208.\\[4pt]
\nolinkurl{pa_tm_research_2026-11-13} & Comparison and reduction research. States a new 196-state variant with shorter source paths and nothing below 196.\\[4pt]
\nolinkurl{pa_consistency_research_2026-11-13} & Comparison and diagonal-vacuity research. States a shorter axiom program with a 199-state implementation and nothing below 196.\\[4pt]
\nolinkurl{conpa_tm_comparison_and_reduction_search} & Comparison and reduction search. States two experimental 196-state command variants and nothing below 196.\\[4pt]
\nolinkurl{pa196_research_2026-11-13} & PA196 comparison, source factorization and proof obligations. States a source tree of 157 internal nodes and nothing below 196.\\[4pt]
\nolinkurl{pa_consistency_research} & Research and verification packet. States that no machine below 196 states has been validated.\\[4pt]
\nolinkurl{pa_tm_comparison_and_reductions} & Comparison and new rewrites, a research snapshot.\\[4pt]
\nolinkurl{pa_tm_extended_research_2026-11-13} & Comparison and refined reduction search. States that the smallest accepted count is 196.\\[4pt]
\nolinkurl{conpa_tm_comparison_and_search} & Comparison and a search below 196. States that none was obtained; files named \file{195} are rejected or diagnostic candidates.\\[4pt]
\nolinkurl{pa_consistency_reduction} & Comparison and reduction experiments. States two new 196-state variants and a shorter source program with a 203-state implementation.\\[4pt]
\nolinkurl{pa_consistency_comparison_2026_11_13} & Comparison and new reduction experiments. States that nothing below 196 was obtained.\\[4pt]
\nolinkurl{PA_consistency_TM_research_2026-11-13} & Comparison and new source reductions. States that no machine below 196 is claimed.\\[4pt]
\nolinkurl{pa_consistency_research_2026-11-13_2} & Comparison and search below 196. States two new layouts with 196 and 198 states and a command-level reduction.\\[4pt]
\nolinkurl{conpa_tm_study_2026-11-13} & Comparison and reduction study. States that nothing below 196 was obtained.\\[4pt]
\nolinkurl{pa_below196_research} & Comparison and reduction experiments. States that nothing below 196 was obtained.\\[4pt]
\nolinkurl{pa_consistency_comparison_and_new196_2026-11-13} & Comparison and reduction research. States that it contains no accepted machine below 196.\\[4pt]
\nolinkurl{pa_tm_research_2026-11-13_2} & Comparison and reduction investigation. States same-size work, failed candidates and nothing below 196.\\[4pt]
\nolinkurl{pa_consistency_cross_archive} & Cross-archive constructions and reduction search. States that it does not certify a 195-state machine.\\[4pt]
\nolinkurl{pa195_conpa_2026-11-13} & PA195. States a 195-state machine \file{pa195/pa195.tm} and an article with a PA proof of the equivalence.\\[4pt]
\nolinkurl{pa_consistency_reduction_audit_2026-11-14} & Comparison and reduction audit. States nothing below 196 and no new theorem.\\[4pt]
\nolinkurl{pa_consistency_research_2026-11-14} & Comparison and NEG18. States an 18-instruction encoded negation with a compiled machine of 197 states.\\[4pt]
\nolinkurl{PA195_ConPA_research_2026-11-14} & A 195-state machine. States the primary \file{machine_195S.tm} and a second table \file{machine_195C.tm}.\\[4pt]
\nolinkurl{pa195_followup_research} & Comparison and searches below 195. States a graph-distinct 195-state machine M195S8 and nothing below 195.\\[4pt]
\nolinkurl{pa194_conpa_2026-11-13} & A 194-state machine. States the table \file{construction/tables/pa194.tm} and a proof of the equivalence in PA.\\[4pt]
\nolinkurl{pa194_conpa_2026-11-13_2} & PA194. States a 194-state machine \file{pa194/pa194.tm} and an article with a PA proof of the equivalence.\\[4pt]
\nolinkurl{pa195_diagonal_combined_2026-11-13} & PA195DC. States a different 195-state machine.\\[4pt]
\nolinkurl{pa194C_conpa_2026-11-13 (2)} & Same PA194C literal table; independent whole-tape symbolic checks, cached command algebra, corpus comparison.\\[4pt]
\nolinkurl{pa194_new_designs_2026-11-13} & Same PA194C plus PA194CS symmetric-Shift implementation; 3646 source sites and additional controller fold.\\[4pt]
\nolinkurl{pa194c_conpa} & Same PA194C; refreshed comparison and secondary PA194altC 68-reader/eight-fold cached search seed.\\[4pt]
\nolinkurl{pa194c_conpa_2026-11-13 (1)} & Same PA194C; supplemental PA194CZ zero-register-jump seed, not a state saving.\\[4pt]
\nolinkurl{pa194c_conpa_2026-11-13} & Same PA194C; inherited calibration, command checks and independent physical checker; 34-archive historical census.\\[4pt]
\bottomrule
\end{longtable}


\section{Independent review and artifact boundary}
A separate mathematical review checked the command proof before packaging.
It confirms the total Pop contract even at formula code one, the one-command
Discard identity for arbitrary incoming cache, the exact ten-entry selectors,
the reserved-zero MP compiler, and the arbitrary-incoming-store refutation
word. It also checks the ordinary proof-annotation converse. The complete
review is supplied as \file{mathematics/independent_review.tex}; it is not
represented as an external human endorsement or a proof-assistant check.

The independent physical review uses different tape representations and
control-flow resolvers, and no compiler or arithmetic-box oracle. Its
portable source and full receipts are under \file{verification/independent/}.
The logical calibration, source-level compiler and physical semantics remain
explicit, separate dependencies in the manifest.

\clearpage
\section*{Sources and attribution}
\addcontentsline{toc}{section}{Sources and attribution}
\begin{enumerate}[leftmargin=*]
\item Vladimir Reshetnikov, \emph{ProveIt}, snapshot
\nolinkurl{17f92105e2bc4317feae6c219cb2143437b482ff}.
\url{https://github.com/VladimirReshetnikov/ProveIt/tree/17f92105e2bc4317feae6c219cb2143437b482ff/Computability/BusyBeaver/Research/PA}.
The comparison ledger identifies the complete packet corpus.
\item Primary separate-head PA194 packet:
\file{PA/docs/pa194_conpa_2026-11-13/}, including the 67-reader/seven-fold
decoder, source map, ordinary calibration and two physical checkers.
The unchanged source fragments and licenses are retained in this archive.
\item Earlier PA command-reduction research, with the no-successful-MP
zero-push interface, and the original ten-entry separate-head calibration.
Exact file identities are recorded in the construction provenance.
\item James Chang's PA196 and PA209 developments, imported under
\file{PA/jamesyc-upstream/}; the retained register compiler, numerical
coding and source-language libraries are credited to their supplied notices.
\item ZFC research, especially
\file{candidate-k215/alternatives/zero-spare-transfer/report.tex}, its
no-successful-MP-zero-push discussion, and the source/physical preservation
methods. \url{https://github.com/VladimirReshetnikov/ProveIt/tree/17f92105e2bc4317feae6c219cb2143437b482ff/Computability/BusyBeaver/Research/ZFC}.
\item ConciseIPRA native-core presentations and their exact interface
reductions. \url{https://github.com/VladimirReshetnikov/ProveIt/tree/17f92105e2bc4317feae6c219cb2143437b482ff/Logic/PrimitiveRecursiveArithmetic/ConciseIPRA}.
No equivalence between $\Con(\mathrm{iPRA})$ and $\Con(\PA)$ is assumed.
\end{enumerate}
\end{document}
